\section{Data and methods}\label{sec:data}

\subsection{Data}

The data are the Danish and Norwegian files of the Luxembourg Income Study (LIS) Database \citep{LIS2024}. Both are drawn from administrative registers. The files of eight further countries, used in Section~\ref{sec:countries}, are described there. The Danish files used here are those for 2004, 2007, 2010, 2013, 2016 and 2019; the first four contain 32{,}000 to 34{,}500 one-person households and the last two about 127{,}000 and 132{,}000. The files from 2015 onwards are about four times as large as the earlier ones, and since tests are more powerful in the larger files, results are compared across that break with care. The composition of the files changes as well. In the later files public contributory pensions are recorded for more older people and at higher amounts: among the transfer-reliant the 90th percentile of these pensions rises from 18{,}000 kroner in 2013 to 144{,}000 in 2016, and a tenth of the group receives less than the full folkepension, presumably reduced pensions for incomplete residence, which do not appear in the earlier files. Section~\ref{sec:results} deals with the first change by a refinement of the transfer-reliant group. The Norwegian files are those for 2007, 2010, 2013, 2016, 2019 and 2022, with 90{,}000 to 113{,}000 one-person households each; the 2004 Norwegian file is a survey with about 3{,}000 and is not used. LIS microdata cannot be taken out of the LIS remote-execution system; every number in this paper is an aggregate computed there.

The unit is the household of one person, identified by the LIS household size variable, and the income concept is LIS disposable household income: the sum of earnings, capital income, pensions and public transfers, less income taxes and social contributions. Because the household has one member, no equivalence scale is needed. Households with non-positive disposable income are dropped. Gross income is the LIS total household income before taxes. The age of the household's only member comes from the LIS person file. Cell shares, population shares and the bootstrap use the LIS household weight; the trimming quantiles, medians and bunching counts are unweighted.

The decomposition by income source uses the LIS harmonised income components: public non-contributory pensions (hi31, which in Denmark contains the folkepension and in the detailed item hi311 its universal part), public contributory pensions (hi32), private pensions (hi33, occupational and individual), public social benefits other than pensions (hipubsoc, which contains housing benefits, hi46), earnings (hilabour) and capital income (hicapital). A person is called \emph{transfer-reliant} if public pensions and public benefits make up at least 90 per cent of the sum of these components. Since public benefits include the housing benefit, receiving it raises a pensioner's chance of being classified as transfer-reliant; this matters only at the margin of the 90 per cent threshold, and the results are reported at other thresholds as well. The threshold is a convention; Section~\ref{sec:results} reports the main test also at 80 and 95 per cent.

\subsection{A test of unimodality for binned incomes}

Let the incomes of a group be binned into cells of width $w$ centred at multiples of $w$, and let $p=(p_1,\dots,p_K)$ be the weighted shares of the cells that lie between two quantiles of the group's income, the $\alpha$-quantile and the 95th percentile. The vector $p$ is \emph{unimodal} if it rises weakly to a maximum and falls weakly afterwards. The departure of $p$ from unimodality is measured by the discrete excess mass of \citet{MullerSawitzki1991}: for a level $\lambda$, let $E_m(\lambda)$ be the largest value of $\sum_{k\in I}(p_k-\lambda)$ over unions $I$ of at most $m$ disjoint runs of consecutive cells, and let
\[
\Delta(p)=\max_{\lambda}\,\{E_2(\lambda)-E_1(\lambda)\}.
\]
Then $\Delta(p)\ge0$, with equality if and only if $p$ is unimodal; $\Delta$ is the discrete counterpart of the excess mass statistic of \citet{MullerSawitzki1991}, which is closely related to the dip of \citet{HartiganHartigan1985}. The test rejects when $\Delta(\hat p)$ is large compared with its distribution under resampling from a unimodal vector close to $\hat p$. That vector is obtained in two steps: adjacent cells whose shares do not differ significantly, at a threshold of $\sqrt{\log n}$ standard errors, are pooled, and the pooled vector is projected onto the unimodal vectors by isotonic regression on either side of each candidate mode, keeping the candidate with the smallest squared distance. The first step matters on the boundary of the null, where $p$ is unimodal but has flat stretches; resampling from the projection alone would treat sampling noise on a flat stretch as a unimodal shape, and the test would be oversized. \citet{Orujov2026} shows that, with a fixed number of cells, the resulting test has asymptotic level at most its nominal level at every unimodal $p$, including those on the boundary of the null, and that it is consistent against every $p$ that is not unimodal. Because the code run on LIS is a compact reimplementation, its size was checked outside LIS by simulation from four unimodal distributions at a sample size of 7{,}000, 100 replications each: a normal and a lognormal distribution, a distribution with a flat top 40{,}000 kroner wide, and a sharp peak with a long right tail. The rejection rates at the five per cent level were 0.00, 0.00, 0.04 and 0.00. The flat-topped distribution, which lies on the boundary of the null, is the hard case, and it was also simulated at larger sizes: at $n=20{,}000$ the rejection rate was 0.055 over 200 replications, and at $n=45{,}000$ it was 0.073 over 300, with a $p$-value of at most 0.005 in two of the 200 replications for which $p$-values were kept. The test is therefore close to its nominal level at the sample sizes used here, slightly liberal at the largest, and a $p$-value of 0.005 is strong evidence at every size. The effective sample size used in resampling is Kish's $(\sum_i v_i)^2/\sum_i v_i^2$ for weights $v_i$. The bootstrap uses 199 draws throughout, with the random seed fixed, so that the smallest attainable $p$-value is $0.005$; resampled vectors are not re-pooled; and $\Delta$ is computed as a maximum over a grid of 1{,}000 levels $\lambda$ together with the observed cell shares.

When the test rejects, the maximising level $\lambda^*$ and the two runs of cells that attain $E_2(\lambda^*)$ locate a valley, the cells strictly between the runs, and two modes, the largest cell to the left of the valley and the largest cell to its right. A $p$-value of one means that the observed cell vector is itself unimodal. These are the locations reported below. They are estimates of the positions of modes and valley at the resolution $w$, and they are reported only when the test rejects.

Two choices need justification. The cell width is 5{,}000 kroner in the main results, a resolution at which a cell of the distribution of older people living alone holds up to eight per cent of the group; the main tests are repeated at 10{,}000 kroner. The lower trimming quantile is $\alpha=0.05$. In every Danish register wave from 2010 onwards a small group has disposable income within a few hundred kroner of zero, often at the same few values, and at $\alpha=0.01$ these records form a spike in the first cell that the test duly reports as a mode. Since that spike is a feature of how the registers record people with essentially no income, not of the distribution of income among people who have one, the main results trim it; the results at $\alpha=0.01$ are in the appendix.

\subsection{Excess mass at a statutory amount}

The concentration of a group's income at a statutory amount $F$ is measured as in the bunching literature \citep{Saez2010,ChettyFriedmanOlsenPistaferri2011,Kleven2016}. Incomes within 40{,}000 kroner of $F$ are counted in bins of 500 kroner; a polynomial of degree five in the distance from $F$ is fitted to the counts outside the window $|x-F|<2{,}500$ and used to predict the counts inside it; and the excess is the difference between observed and predicted counts in the window, relative to the predicted count. Nothing in the present application calls for the structural interpretation that the bunching literature gives this number. It is used only as a descriptive measure of whether the distribution has more mass at $F$ than its shape nearby would imply, and its reliability is checked by computing the same quantity at the placebo amounts $0.8F$, $0.9F$, $1.1F$ and $1.2F$.
